{\large\textbf{Sand craters and dunes (10.0 points)}}

NASA's Spirit rover (\textbf{Fig.~1.}(a)) landed on Mars in 2004 to study its
geology and potential presence of water. The landing site (\textbf{Fig.~1.}(b))
is surrounded by craters of various sizes and sand dunes. During exploration,
the rover must avoid getting stuck in the sand dunes of Mars.

\begin{center}\includegraphics[width=0.81\linewidth]{figures/IPhO_2025_Q5_fig1.jpg}\end{center}

\begin{center}\textbf{Fig.~1.} (a) Artist's view of Spirit. (b) Landing site of
the rover on Mars. The scale bar represents $200\,\mathrm{m}$.\end{center}

The problem has two independent parts A (crater formation) and B (sand
trapping) that can be treated in any order. The list of equipment is given
below and illustrated in \textbf{Fig.~2.}
\begin{itemize}
\item (a) Plastic box, needs to be emptied. The empty box will be used to
collect the overflowing sand during experiments.
\item (b) Bowl.
\item (c) Bottle of sand.
\item (d) 6 steel balls in a container. The balls have 4 different diameters.
The three smallest ones are identical.
\item (e) Tape measure.
\item (f) Holding device consisting of a wooden tray with rubber feet (f1), a
vertical rod (f4), clamping screw (f2) and horizontal rod (f3). The different
elements must be assembled as shown in the photo (f).
\item (g) Sieve, used to find the small ball if it gets lost in the sand.
\item (h) Aluminium rail, $1\,\mathrm{m}$ long.
\item (i) Brush to clean the rail and balls of sand if necessary.
\item (j) Wooden track.
\item (k) Chronometer.
\item (l) Adhesive putty.
\item (m) Funnel to help to put the sand back into the box at the end.
\item (n) Spoon.
\item (o) Ruler.
\end{itemize}

\begin{center}\includegraphics[width=0.77\linewidth]{figures/IPhO_2025_Q5_fig2.jpg}\end{center}

\begin{center}\textbf{Fig.~2.} Photographs of all equipment.\end{center}

\textbf{A. Impact craters}

Craters on Mars, whose diameter $D$ varies from about $10\,\mathrm{m}$ to
several hundreds of km, result from the impact of meteorites. Different models
predict how $D$ depends on the impact parameters: impactor diameter $d$, energy
$E$ (\textbf{Fig.~3}).

\begin{center}\includegraphics[width=0.72\linewidth]{figures/IPhO_2025_Q5_fig3.png}\end{center}

\begin{center}\textbf{Fig.~3.} Crater formation.\end{center}

\textbf{Model 1:} $D$ depends only on the impactor diameter $d$
\begin{equation}
D = c_1 d, \tag{1}
\end{equation}
where $c_1$ is a dimensionless number independent of $E$ and $d$.

\textbf{Model 2:} the meteorite energy $E$ is converted through volumic
processes during the impact. This model predicts that $D$ is proportional to
$E^{1/3}$
\begin{equation}
D = c_2 E^{1/3} \tag{2}
\end{equation}
where $c_2$ is a parameter independent of $E$ and $d$.

\textbf{Model 3:} $E$ is used to eject material outside the crater. Under this
assumption
\begin{equation}
D = c_3 E^{1/4} \tag{3}
\end{equation}
where $c_3$ is a parameter independent of $E$ and $d$.

Here, we perform experiments on crater formation at a centimeter scale to
compare the three models. Steel balls of different diameters $d$ and masses
$m$, with a density
$\rho_a = 7.8 \times 10^3\,\mathrm{kg \cdot m^{-3}}$ (item (\textbf{d}) of the
equipment list), act as the meteorites.

\begin{center}
\begin{tabular}{|l|l|l|}
\hline
Ball \#1 & $d_1 = 2.0\,\mathrm{mm}$ & $m_1 = 0.033\,\mathrm{g}$ \\ \hline
Ball \#2 & $d_2 = 5.0\,\mathrm{mm}$ & $m_2 = 0.51\,\mathrm{g}$ \\ \hline
Ball \#3 & $d_3 = 9.0\,\mathrm{mm}$ & $m_3 = 3.0\,\mathrm{g}$ \\ \hline
Ball \#4 & $d_4 = 16.0\,\mathrm{mm}$ & $m_4 = 17\,\mathrm{g}$ \\ \hline
\end{tabular}
\end{center}

The bowl (\textbf{b}) filled with sand (\textbf{c}) is placed inside the emptied
plastic box (\textbf{a}) that will help collect the excess sand. The bowl is
filled completely with sand and the surface is carefully leveled with the edge
of the ruler (\textbf{o}). Avoid compacting the sand! To release the ball above
the bowl, one can use the stand equipped with a rod and thumbscrew
(\textbf{f}). The rod serves as a guide to release the ball directly above the
bowl and also to measure the drop height $h$ above the surface, which will be
measured using the tape measure (\textbf{e}).

\begin{center}\includegraphics[width=0.62\linewidth]{figures/IPhO_2025_Q5_fig4.png}\end{center}

\begin{center}\textbf{Fig.~4.} Crater formation experimental setup.\end{center}

Drop ball \#3 from a height $h = 50\,\mathrm{cm}$ and measure the diameter $D$
of the crater formed. Repeat the experiment 5 times. After each impact, mix the
sand with the spoon (\textbf{n}), and level it carefully with the edge of ruler
(\textbf{o}). Avoid compacting the sand! If needed, use the sieve (\textbf{g})
to find the ball if it gets lost in the sand.

\textbf{A.1}\quad Present your results in a table and give $D$ with its
uncertainty. \hfill 0.6pt

During the fall, the air drag force is
\begin{equation}
F = \frac{1}{8}\pi d^2 \rho_0 C_x v^2 \tag{4}
\end{equation}
where $v$ is the ball velocity,
$\rho_0 \simeq 1.2\,\mathrm{kg \cdot m^{-3}}$ is the air density and $C_x$ is
a dimensionless coefficient of order unity.

The air drag force is negligible if the ball is dropped from a height less than
the maximum drop height $h_\mathrm{max}$, defined as the height at which the air
drag force remains less than 10 \% of the weight throughout the fall.

\textbf{A.2}\quad Determine the theoretical expression for the maximum drop
height $h_\mathrm{max}$. Calculate $h_\mathrm{max}$ numerically for the four
available balls. \hfill 0.5pt

Investigate the relationship between $D$ and $E$ experimentally in order to
compare the three power laws presented in the introduction. Find out if the
exponent changes across the range of energies tested. To achieve this, take a
series of measurements by dropping the balls from different heights. A wide
range of energies must be covered. The balls can be dropped from heights of up
to $h = 2\,\mathrm{m}$ in order to reach high values of $E$ while respecting the
condition established in \textbf{A.2.} For each set of parameters, repeat the
experiment only twice, and compute the mean value $D$.

\textbf{A.3}\quad Present your results in a table: mass of the ball $m$, drop
height $h$, impact energy $E$, crater diameter $D$. \hfill 1.7pt

\textbf{A.4}\quad Plot your results on the graph paper of your choice
(logarithmic or linear). On the graph representation, add lines corresponding
to models 1, 2 and 3. State which of the three theoretical models best fits the
experimental data. \hfill 1.2pt

\textbf{B. Rolling and bogging in sand}

Five years after landing, the rover Spirit bogs in the sands of a Martian dune
for good. Rolling in sand is particularly delicate as the motion of grains
dissipates a lot of energy. Here, we study the braking of a ball rolling in
sand. The ball, initially at rest, is first accelerated on a rail inclined at an
angle $\theta$, then slowed down on a bed of sand.

\begin{center}\includegraphics[width=0.82\linewidth]{figures/IPhO_2025_Q5_fig5.png}\end{center}

\begin{center}\textbf{Fig.~5.} Inclined rail (\textbf{h}) combined with the
wooden track (\textbf{j}).\end{center}

\textbf{Ball motion along the rail}

Ball \#4 is released with no initial speed from an arbitrary point on the rail
(\textbf{h}), chosen as the origin of the $x$-axis ( $x = 0$ )
(\textbf{Fig.~5}). Let $x(t)$ denote the position of the ball along the rail.
The moment of inertia of a ball of mass $m$ and diameter $d$ with respect to an
axis passing through it center is given by $J = md^2/10$. The kinetic energy $K$
of a ball moving at speed $v$ while rotating at angular speed $\omega$ is
\begin{equation}
K = \frac{1}{2}mv^2 + \frac{1}{2}J\omega^2. \tag{5}
\end{equation}

We assume that the ball rolls on the rail without slipping and neglect any
energy dissipation.

\textbf{B.1}\quad Express the position $x$ of the ball as a function of time
$t$, angle $\theta$ and acceleration of gravity $g$. \hfill 0.4pt

One end of the rail (\textbf{h}) rests on the edge of the wooden track
(\textbf{j}), which is at this point empty of sand. The other end of the rail is
supported by the stand (\textbf{f}) in such a way that it forms an angle of
inclination $\theta = 5^\circ$ with the horizontal. Make sure to perform this
adjustment carefully. The rail is secured in place (on both sides) using
adhesive putty (\textbf{l}).

Use a chronometer (\textbf{k}) to measure the time $t_{50}$ taken by the ball
to travel a distance $l = 50\,\mathrm{cm}$ along the rail.

\textbf{B.2}\quad Take 5 measurements and present the result along with the
order of magnitude of its statistical uncertainty. \hfill 0.7pt

Measure $t$ with the order of magnitude of its statistical uncertainty for at
least 8 different values of $\ell$.

\textbf{B.3}\quad Present your results in a table. \hfill 0.8pt

\textbf{B.4}\quad Plot your results with error bars to confirm the law
established at question \textbf{B.1}. Deduce an experimental estimate of the
constant $g$ with its uncertainty. \hfill 1pt

\textbf{Motion of the ball in sand}

We note $\ell$ the distance travelled by the ball on the rail. On the sand, the
ball comes to a stop after travelling a distance $L$ as defined in
\textbf{Fig.~6}.

\begin{center}\includegraphics[width=0.82\linewidth]{figures/IPhO_2025_Q5_fig6.png}\end{center}

\begin{center}\textbf{Fig.~6.} Acceleration over a distance $\ell$ and stopping
over a distance $L$.\end{center}

It is thus slowed down by a drag force $T$ which may have two possible origins:
\begin{itemize}
\item \textbf{Model \#1 (solid friction):} as between two solids in relative
motion, the sand exerts on the ball a constant drag force
$T = -\mu_\mathrm{eff} mg$ , where $\mu_\mathrm{eff}$ is the effective drag
coefficient of the ball-sand contact and $m$ is the mass of the ball.
\item \textbf{Model \#2 (fluid drag):} the drag force depends linearly on the
ball velocity, $T = -kv$ where $k$ is a constant and $v$ the norm of the
velocity.
\end{itemize}

The goal here is to determine which proposition best describes the observed
braking behavior.

When moving in sand, the ball is modelled as a point mass. Given the small
value of the slope of the rail, we neglect any energy loss in the transition
between the rail and the sand track. Establish the theoretical law linking $L$
to $\ell$ in each of the two situations (solid friction or fluid drag). The two
suggestions lead to a power law of the form $L \sim \ell^\alpha$ in which the
exponent $\alpha$ takes two different values.

\textbf{B.5}\quad For model 1 and model 2, give the relationship between $L$
and $\ell$ and the value of $\alpha$. \hfill 0.6pt

Place the wooden track (\textbf{j}) on a sheet of paper. Fill the track with
sand and prepare a uniform layer by carefully scraping the surface with the
ruler. Avoid compacting the sand! Adjust again carefully the angle of the rail
to $\theta = 5^\circ$. Release ball \#4 ($d_4 = 16.0\,\mathrm{mm}$) on the
inclined rail so that the distance travelled on the rail is
$l = 50\,\mathrm{cm}$.

\textit{Before each run, stir the sand, refill the track and scrape the surface
again. Clean the rail and the ball from sand by using the brush
(\textbf{i}). At the end of the experiment, use the sheet of paper as a funnel
to put the sand in excess back in the bottle.}

\textbf{B.6}\quad Measure the distance $L_{50}$ travelled in the sand until the
ball comes to a stop. Perform several measurements (at least 5) to determine
$L_{50}$ along with its unit and uncertainty. \hfill 0.8pt

\textbf{B.7}\quad After several measurements for at least 8 values of $\ell$
(keeping $\theta = 5^\circ$), plot $L$ with its error bars as a function of
$\ell$ and conclude which model best describes the drag force $T$. \hfill 1.5pt

\textbf{B.8}\quad Based on the chosen model, specify the value of the
coefficient $\mu_\mathrm{eff}$ or $k$ that characterizes the force $T$.
\hfill 0.2pt
